% part: intuitionistic-logic
% chapter: propositions-as-types
% section: sequent-natural-deduction

\documentclass[../../../include/open-logic-section]{subfiles}

\begin{document}

\olfileid{int}{pty}{snd}

\olsection{क्रमवर्ती-रूप नैसर्गिक निगमन}

$\Gamma$ ही !!{undischarged} गृहीतके आणि~$!A$ हा
निष्कर्ष असलेली नैसर्गिक निगमनाची !!{derivation} असेल,
तर $\Gamma \Sequent !A$ असे लिहू;~$\Gamma$ रिकामे
असेल, तर $\Sequent !A$ असे लिहू.

$\Gamma \cup \{!A_1, \dots, !A_n\}$ साठी
$\Gamma, !A_1, \dots, !A_n$ आणि $\Gamma \cup \Delta$
साठी $\Gamma, \Delta$ असे लिहू.

$\Gamma \Sequent !A \land !B$ असेल, म्हणजे~$\Gamma$
ही !!{undischarged} गृहीतके आणि $!A \land !B$ हे
अंतिम !!{formula} असलेली !!a{derivation} असेल,
तर तळाशी $\Elim{\land}$ लावून तीच
!!{undischarged} गृहीतके आणि~$!A$ हा निष्कर्ष
असलेली !!a{derivation} मिळते. दुसऱ्या शब्दांत,
$\Gamma \Sequent !A \land !B$ असेल, तर
$\Gamma \Sequent !A$.
\begin{prooftree}
  \Axiom$\Gamma \fCenter !A \land !B$
  \RightLabel{$\Elim{\land}$}
  \UnaryInf$\Gamma \fCenter !A$
  \DisplayProof\qquad\bottomAlignProof
  \Axiom$\Gamma \fCenter !A \land !B$
  \RightLabel{$\Elim{\land}$}
  \UnaryInf$\Gamma \fCenter !B$
\end{prooftree}
$\Elim{\land}$ ही खूण नैसर्गिक निगमनातील त्याच
नावाच्या नियमाशी असलेला संबंध सूचित करते.

त्याचप्रमाणे $\Gamma, !A \Sequent !B$ असेल,
म्हणजे $\Gamma, !A$ ही !!{undischarged} गृहीतके
आणि~$!B$ हे अंतिम !!{formula} असलेली
!!a{derivation} असेल, तर $\Intro{\lif}$ नियम लावून
$\Gamma$ ही !!{undischarged} गृहीतके आणि
$!A \lif !B$ हे अंतिम !!{formula} असलेली
!!a{derivation} मिळते; म्हणजे
$\Gamma \Sequent !A \lif !B$.
यात गृहीतकांचे !!{discharge} अधिक स्पष्ट झाले आहे.
\begin{prooftree}
  \Axiom $\Gamma, !A \fCenter !B$
  \RightLabel{$\Intro{\lif}$}
  \UnaryInf $\Gamma \fCenter !A \lif !B$
\end{prooftree}

इतर नियमांवरूनही अशाच प्रकारे निष्कर्ष काढता
येतात. ते पुढीलप्रमाणे स्पष्ट केले आहेत:
\begin{gather*}
  \Axiom $\Gamma \fCenter !A$
  \Axiom $\Delta \fCenter !B$
  \RightLabel{$\Intro{\land}$}
  \BinaryInf $\Gamma,\Delta \fCenter !A \land !B$
  \DisplayProof\\
  \Axiom $\Gamma \fCenter !A \land !B$
  \RightLabel{$\Elim{\land}_1$}
  \UnaryInf $\Gamma \fCenter !A$
  \DisplayProof
  \qquad
  \Axiom $\Gamma \fCenter !A \land !B$
  \RightLabel{$\Elim{\land}_2$}
  \UnaryInf $\Gamma \fCenter !B$
  \DisplayProof
  \\
  \Axiom $\Gamma \fCenter !A$
  \RightLabel{$\Intro{\lor}_1$}
  \UnaryInf $\Gamma \fCenter !A \lor !B$
  \DisplayProof\qquad
  \Axiom $\Gamma \fCenter !B$
  \RightLabel{$\Intro{\lor}_2$}
  \UnaryInf $\Gamma \fCenter !A \lor !B$
  \DisplayProof\\
  \Axiom $\Gamma \fCenter !A \lor !B$
  \Axiom $\Delta, !A \fCenter !C$
  \Axiom $\Delta', !B \fCenter !C$
  \RightLabel{$\Elim{\lor}$}
  \TrinaryInf $\Gamma, \Delta, \Delta' \fCenter !C$
  \DisplayProof
  \\
  \Axiom $\Gamma, !A \fCenter !B$
  \RightLabel{$\Intro{\lif}$}
  \UnaryInf $\Gamma \fCenter !A \lif !B$
  \DisplayProof
  \qquad
  \Axiom $\Delta \fCenter !A \lif !B$
  \Axiom $\Gamma \fCenter !A$
  \RightLabel{$\Elim{\lif}$}
  \BinaryInf $\Gamma, \Delta \fCenter !B$
  \DisplayProof
  \\
  \Axiom $\Gamma \fCenter \lfalse$
  \RightLabel{$\FalseInt$}
  \UnaryInf $\Gamma \fCenter !A$
  \DisplayProof
\end{gather*}

कोणतेही गृहीतक स्वतःच~$!A$ वरून~$!A$ ची
!!a{derivation} असते; म्हणजे $!A \Sequent !A$
हे नेहमीच असते.

\begin{prooftree}
  \AxiomC{$ $}
  \UnaryInfC{$!A \fCenter !A$}
\end{prooftree}

हे नियम एकत्रितपणे कोणत्या नैसर्गिक निगमनाच्या
!!{derivation}s अस्तित्वात आहेत यांचे कलन म्हणून
घेता येतात. त्यांना नैसर्गिक निगमनाचे चिन्हांकनात्मक
रूपांतरही मानता येते; प्रत्येक टप्प्यावर कोणते
!!{formula} !!{derive} झाले आहे एवढेच नव्हे,
तर ते कोणत्या !!{undischarged} गृहीतकांवरून
!!{derive} झाले आहे हेही त्यात नोंदते.

\begin{prooftree}
  \Axiom$!A \fCenter !A$
  \UnaryInf $!A \fCenter !A \lor (!A \lif \lfalse)$
  \Axiom $!B \fCenter !B$
  \BinaryInf$!A, !B\lif \fCenter \lfalse$
  \UnaryInf$(!B \fCenter !A \lif \lfalse$
  \UnaryInf$( !B \fCenter !A \lor (!A \lif \lfalse)$
  \Axiom$(!B \fCenter !B$
  \BinaryInf$(!B \fCenter \lfalse$
  \UnaryInf$\fCenter !B \lif \lfalse$
\end{prooftree}
येथे~$!B$ हे $(!A \lor (!A \lif \lfalse))\lif \lfalse$
याचे संक्षिप्त रूप आहे.

\end{document}
